\chapter{Introduction}
\label{ch:intro}

Static code analysis has become a fundamental technique in modern software engineering, allowing developers to identify potential defects, vulnerabilities, and code quality issues without executing the program. Unlike dynamic analysis methods that require program execution, static analysis examines source code structure, control flow, and data flow to detect anomalies that could lead to runtime errors, security vulnerabilities, or maintenance difficulties. The effectiveness of static analysis tools has made them essential components of continuous integration pipelines and quality assurance processes in software development organizations.

CodeChecker is a static analysis infrastructure built on the LLVM/Clang Static Analyzer toolchain, designed to replace the traditional \texttt{scan-build} tool in Linux and macOS development environments. The system provides a comprehensive platform for managing static analysis results from multiple analyzers, including Clang Static Analyzer, Clang Tidy, Cppcheck, and various other tools supporting C/C++, Java, Python, and JavaScript. CodeChecker offers both command-line interfaces and web-based visualization tools, enabling teams to store, query, filter, and compare analysis results across multiple code revisions and analysis runs.

While static analysis tools excel at identifying potential code defects, they operate independently of test execution and code coverage metrics. Test coverage analysis measures which lines of source code are executed during test runs, providing developers with insights into the extent to which their test suites exercise the codebase. The integration of code coverage visualization with static analysis results presents a valuable opportunity to correlate static analysis findings with actual test coverage, enabling developers to identify areas where both static analysis warnings exist and test coverage is insufficient.

This thesis presents the design, implementation, and evaluation of a code coverage visualization feature integrated into the CodeChecker static analysis platform. The feature extends CodeChecker's capabilities by allowing users to generate test coverage data, store it alongside static analysis results, and visualize coverage information through an interactive web interface that highlights covered and uncovered lines of source code. This integration enables developers to make more informed decisions about code quality by simultaneously considering both static analysis findings and test coverage metrics.

The implementation required knowledge beyond the compulsory curriculum, including expertise in static analysis toolchains, database schema design, web application frameworks, and coverage instrumentation techniques. The solution integrates multiple platforms and technologies: command-line tools for coverage generation, database systems for data persistence, and web frameworks for visualization. The feature is part of the larger CodeChecker ecosystem, requiring deep understanding of its architecture, API design patterns, and database models before implementation could begin. The implementation involved non-trivial algorithms for parsing standardized coverage data formats (LCOV), efficient database query design for large-scale coverage data, and performance optimization for real-time code visualization in web browsers.

The remainder of this thesis is structured as follows. Chapter~\ref{ch:user} provides user documentation describing how to use the code coverage feature, including command-line tools for generating coverage data and the web interface for visualization. Chapter~\ref{ch:impl} presents the technical implementation details, including the database schema design, API endpoints, and frontend components. Finally, Chapter~\ref{ch:sum} summarizes the contributions and discusses potential future enhancements to the feature. 
